\section{The exact task defect at a prescribed output width}\label{sec:width-defect}
The raw matrix order of Section~\ref{sec:order} is exact for localized affine-square tests. Ordinary squared prediction tests form a smaller family. Requiring the stronger order can therefore discard valid task comparisons, while taking a supremum over arbitrarily many readouts can overcharge a finite decoder. We identify the least additive defect for the actual task and its output width.

\subsection{Finite dictionaries and an attained optimal defect}
Let $Y,Z$ be standard Borel spaces, $D$ a nonempty compact metric decision space, and $\ell:Z\times D\to[0,L]$ Borel. Assume a uniform decision modulus
\begin{equation}\label{eq:loss-modulus}
 |\ell(z,r)-\ell(z,s)|\le\omega(d(r,s)),\qquad
 \lim_{\delta\downarrow0}\omega(\delta)=0,
\end{equation}
with $\omega$ nondecreasing. For finite positive measures $P,Q$ on $Y\times Z$, put $\lambda=P_Y+Q_Y$. The unnormalized conditional loss densities $g_P(y,r),g_Q(y,r)$ relative to $\lambda$ have jointly Borel versions continuous in $r$. Indeed disintegrate $P+Q$ over $\lambda$, multiply its conditional law by the Radon--Nikodym densities of $P,Q$ relative to $P+Q$, and integrate $\ell$. Off one completed null set,
\[
 |g_P(y,r)-g_P(y,s)|\le p(y)\omega(d(r,s)),\quad
 |g_Q(y,r)-g_Q(y,s)|\le q(y)\omega(d(r,s)),
\]
where $p=dP_Y/d\lambda$, $q=dQ_Y/d\lambda$ and $p+q=1$. Zero total mass causes no difficulty. Put $F(y,r)=g_P(y,r)-g_Q(y,r)$ and
\[
 \Delta_C^\ell(P\mid Q)=\int\max\{0,\max_{r\in C}F(y,r)\}\,d\lambda
\]
for a nonempty finite dictionary $C\subset D$.

\begin{theorem}[Width-resolved optimal task defect]\label{thm:width-defect}
For every integer $m\ge1$, the following equal quantities are attained:
\begin{align}
 \Delta_m^\ell(P\mid Q)&=\max_{(r_1,\ldots,r_m)\in D^m}
 \int\max\{0,F(y,r_1),\ldots,F(y,r_m)\}\,d\lambda,
 \label{eq:defect-envelope}\\
 &=\sup_{\substack{a:Y\to[0,1]\ {
m Borel}\\r:Y\to D\ {
m Borel},\ |r(Y)|\le m}}
       \int a(y)\ell(z,r(y))\,(P-Q)(dy,dz).
 \label{eq:defect-tests}
\end{align}
Thus this is the smallest additive error valid for every selected width-$m$ task. An optimum uses a deterministic report event and readout from one fixed dictionary; independent current randomization cannot improve it. The value is independent of the dominating measure and satisfies
\begin{equation}\label{eq:width-growth}
 \Delta_1\le\Delta_m\le m\Delta_1,\quad
 \Delta_{m+n}\le\Delta_m+\Delta_n,\quad
 \frac{\Delta_n}{n}\le\frac{\Delta_m}{m}\quad(1\le m\le n).
\end{equation}
Moreover
\begin{equation}\label{eq:infinite-defect}
 \Delta_\infty^\ell(P\mid Q):=\int[\sup_{r\in D}F(y,r)]_+\,d\lambda
       =\lim_{m\to\infty}\Delta_m^\ell(P\mid Q).
\end{equation}
If a $\delta$-net of $D$ has $n$ points, then
\begin{equation}\label{eq:defect-net}
 0\le\Delta_\infty-\Delta_n\le
 \omega(\delta)[P(Y\times Z)+Q(Y\times Z)].
\end{equation}
In particular $\Delta_1=0$ if and only if $\Delta_\infty=0$. The width hierarchy distinguishes positive defect magnitudes, not different zero-defect orders.
\end{theorem}
\begin{proof}
Fix an ordered dictionary. At each report choose the first index maximizing its positive loss-density difference, and accept exactly when that maximum is positive. These decisions are Borel and give the displayed envelope. They pointwise dominate every other selection and every randomized readout in that dictionary. Every finite-image rule uses some such dictionary, proving equality of the suprema.

Moving every tuple coordinate by at most $\delta$ changes the envelope by at most $(p+q)\omega(\delta)$ almost everywhere. Integration makes the objective continuous on compact $D^m$, with repetitions allowed. It attains its maximum, and the selected deterministic test attains the second supremum. Averaging over an independent current seed cannot increase that value. A seed choosing dictionaries whose union exceeds $m$ does not describe a width-$m$ decoder.

The test expression gives independence of the dominating measure. Inclusion gives monotonicity. A nonnegative maximum is bounded by the sum of the single-center positive parts, proving $\Delta_m\le m\Delta_1$. Partition a dictionary into two groups to obtain subadditivity. For $m\le n$, choose an optimal $n$-tuple and average over its $m$-element index subsets. A selected maximizing positive index is retained with probability $m/n$. Thus the average subdictionary envelope is at least $m/n$ times the original one. Integrate and take one subdictionary at least as good as its average. This proves the last inequality, including repeated entries.

Continuity makes a countable dense decision set determine the pointwise supremum. The positive envelopes of its finite initial segments increase to that supremum; monotone convergence proves the limit and measurability. A finite net gives the quantitative estimate by the same modulus calculation. Finally, $\Delta_1=0$ makes $F(y,r)\le0$ almost everywhere for each point in the countable dense set. Intersect these full-measure sets and use continuity to obtain the inequality for all decisions on one common set. Hence $\Delta_\infty=0$; the converse is immediate.
\end{proof}

For fixed $P_0,P_1$, the map $t\mapsto\Delta_m^\ell(tP_0\mid P_1)$ is convex, nondecreasing and $LP_0(Y\times Z)$-Lipschitz on $[0,\infty)$. It is a supremum of affine functions of $t$ with slopes between zero and that bound. The positive-variation and measurable-selection ingredients are classical. The decision-cardinality restriction should not be confused with Blackwell $k$-decision orders \cite{BR14}: the present localized fixed-loss order at zero defect already follows from width one.

\subsection{The exact squared-task order}
Take $D=H\subset\R^k$ compact convex and $\ell(z,r)=\|z-r\|^2$, with bounded audits. Relative to a common report measure, write the signed moment densities of $P_1-tP_0$ as
\[
 a=\frac{d(P_1-tP_0)_Y}{d\lambda},\quad
 b=\frac{d\int z(P_1-tP_0)(dy,dz)}{d\lambda},\quad
 c=\frac{d\int\|z\|^2(P_1-tP_0)(dy,dz)}{d\lambda}.
\]
The corresponding quadratic is $\psi_y(r)=a(y)\|r\|^2-2b(y)\cdot r+c(y)$.

\begin{theorem}[Necessary and sufficient squared-task certificate]\label{thm:quadratic-defect}
The least unrestricted-width defect is
\begin{equation}\label{eq:quadratic-defect}
 \Delta_\infty^{\rm sq}(tP_0\mid P_1)
     =\int[-\min_{r\in H}\psi_y(r)]_+\,d\lambda.
\end{equation}
Localized squared-task order with zero defect holds for all finite dictionaries if and only if $\min_H\psi_y\ge0$ on one common full-measure set. No inverse covariance, equal matrix support or positive eigenvalue floor is required. For $a>0$ the minimum is
\begin{equation}\label{eq:quadratic-projection}
 c-\|b\|^2/a+a\,\dist(b/a,H)^2;
\end{equation}
for $a=0$ it is $c-2\sup_H b\cdot r$. On the radius-$R$ closed Euclidean ball,
\begin{equation}\label{eq:quadratic-ball}
 \min_H\psi=\begin{cases}
 c-\|b\|^2/a,&a>0,\ \|b\|\le aR,\\
 c+aR^2-2R\|b\|,&\text{otherwise}.
 \end{cases}
\end{equation}
The second line includes the nonpositive-mass and zero-radius boundary cases whenever the first does not apply.
\end{theorem}
\begin{proof}
Expanding the square makes the directed loss-density difference for $(tP_0,P_1)$ equal to $-\psi$. The preceding theorem proves the integral and common-null-set statements. Completing the square proves the first formula; $a=0$ gives a linear optimization. On a ball with $a>0$, either the unconstrained minimizer $b/a$ lies inside or its radial projection minimizes. If $a\le0$, the least value at radius $s$ is $c+as^2-2\|b\|s$, nonincreasing for $0\le s\le R$. Its minimum is at $R$. These arguments include ties, $b=0$ and $R=0$.
\end{proof}

Report mass is a separate resource test: its least selected-event defect in $P_Y\le uQ_Y+e$ is $\int[p-uq]_+d\lambda$. It is not implied by one arbitrary task inequality. Equality of the entire task family has a sharper conclusion.

\begin{proposition}[Task equality and its exact moment content]\label{prop:task-rigidity}
Suppose the affine hull of $H$ has dimension at least one. Localized squared-task equality for $P,Q$ at every constant readout in $H$ is equivalent to equality of their report masses, their first-moment measures projected onto the direction space of that affine hull, and their full centered squared-norm measures $\int\|z-h_0\|^2P(dy,dz)$ for an origin $h_0$ in the affine hull. Full second-moment matrix measures need not agree. For a singleton carrier the mass must instead be tested separately.
\end{proposition}
\begin{proof}
Choose a countable dense set of readouts and intersect their full-measure equality sets. The common modulus extends equality to all of $H$. Translate by $h_0$ and restrict the quadratic difference to the direction space $V$ of the affine hull. A compact convex set with this affine hull has nonempty relative interior. A polynomial vanishing there vanishes on $V$: restrict successively to line segments in a relative open ball, where a one-variable polynomial with an interval of zeros is identically zero, and identify its coefficients. The quadratic coefficient is the scalar mass difference times the identity on $V$, so that difference is zero because $\dim V\ge1$. The linear coefficient is the projected centered first-moment difference, and the constant coefficient is the full centered squared-norm difference. They vanish. Conversely these three equalities make every expanded squared-task difference zero. When $H$ is a singleton, only one constant loss measure is tested, so a zero-loss audit does not identify mass.
\end{proof}

\begin{example}[Identical tasks with incompatible matrix supports]\label{ex:strict-task-order}
There is one report and the decision carrier is the unit ball in $\R^2$. Let $P_0$ audit the two signs of $e_1$ fairly and $P_1$ the two signs of $e_2$ fairly. Every ordinary squared decision has loss $1+\|r\|^2$ under both laws. Their affine-square moment matrices are respectively $\operatorname{diag}(1,1,0)$ and $\operatorname{diag}(1,0,1)$ in the coordinates $(1,z_1,z_2)$. No two-sided positive relative matrix order can compare their different kernels. Thus exact squared-task comparison is strictly broader than the sufficient positive-matrix criterion, without contradicting that criterion's exact statement for its larger test family.
\end{example}

\begin{theorem}[Sharp width amplification in every finite dimension]\label{thm:sharp-width}
For every integer $k\ge1$ there are probability report--audit laws with identical report masses and squared decisions in the unit ball of $\R^k$ such that
\begin{equation}\label{eq:sharp-width-law}
 \Delta_m^{\rm sq}(P_0\mid P_1)=\frac{\min(m,k)}{5k}
 \quad(m\ge1).
\end{equation}
Both the linear width growth and the discrepancy between fixed-width and unrestricted defects are sharp in this family.
\end{theorem}
\begin{proof}
Let the report $i$ be uniform on $\{1,\ldots,k\}$. Under $P_0$ the audit is zero. Under $P_1$, conditional on $i$, it is $e_i$ with probability $4/5$ and $-e_i$ with probability $1/5$. The conditional loss difference is $(6/5)r_i-1$. It is positive only when $r_i>5/6$. A unit-ball vector cannot meet this condition for two different coordinates, since $2(5/6)^2>1$. Its positive contribution in its one possible coordinate is at most $1/5$. A dictionary of $m$ vectors therefore contributes at most $\min(m,k)/(5k)$. Choosing that many distinct basis vectors and accepting their corresponding reports attains the bound. The construction has bounded audits, exact report-mass equality, and no asymptotic dimension assumption.
\end{proof}

\subsection{The same dictionary across actual histories}
Let $u$ range over a standard Borel context space with finite positive occupation measure $\mu$. Let $P_u,Q_u$ be finite kernels with integrable total masses; assume the same bounded decision-continuous loss. Form the finite measures $\mathbf P=\mu(du)P_u(dy,dz)$ and $\mathbf Q=\mu(du)Q_u(dy,dz)$, where the full report is $(u,y)$.

\begin{proposition}[Exact fixed-dictionary occupation defect]\label{prop:occupation-defect}
One has
\begin{equation}\label{eq:occupation-defect}
 \Delta_m^\ell(\mathbf P\mid\mathbf Q)
   =\max_{C\in D^m}\int\Delta_C^\ell(P_u\mid Q_u)\,\mu(du)
   \le\int\Delta_m^\ell(P_u\mid Q_u)\,\mu(du).
\end{equation}
The maximum uses a single dictionary for every context. The inequality can be strict by an arbitrarily large factor as the number of contexts grows. Neither side asserts that a prescribed occupation measure is causally realizable.
\end{proposition}
\begin{proof}
For a fixed dictionary, take the positive envelope of its unnormalized loss densities on the product context/report space. Disintegration and integration give exactly the integral of the rowwise fixed-dictionary envelopes. This can also be checked without selecting contextwise densities: a finite partition of reports and rational acceptance/readout choices approximate every bounded selected task in the finite measure $P_u+Q_u$. A countable algebra generating the standard Borel report space, together with a countable dense decision set, suffices for these approximations. Hence the rowwise suprema are measurable. The uniform decision modulus and integrable total masses make the integrated fixed-dictionary envelope continuous on compact $D^m$, so it attains its maximum. Applying the width-defect theorem to the product measures yields the equality. Moving the maximum inside the integral gives the inequality.

For strictness use the preceding sharp example with context $u=i$ and no additional report. Every row has one-center defect $1/5$, whereas a single common width-$m$ dictionary over uniform contexts has defect $\min(m,k)/(5k)$. Thus even individually exact rowwise errors may overcharge a common finite decoder. In a stopped experiment $\mu$ must be the actual pre-call occupation of the competitor; this algebraic identity does not optimize or fabricate that occupation.
\end{proof}
